% Pollis New Methods: GLM inference (parametric bootstrap and delta-beta influence) — draft 1 (11 October 2026)
% Compile with: latexmk -pdf nm_glm_inference.tex && latexmk -c
% Shared look for the Pollis New Methods documents (same fonts and colours as the brochure and the workflow handouts).
\documentclass[10pt,a4paper]{article}

\usepackage[a4paper,left=18mm,right=18mm,top=15mm,bottom=17mm,footskip=8mm]{geometry}
\usepackage[T1]{fontenc}
\usepackage[default,semibold]{sourcesanspro}
\usepackage{textcomp}
\usepackage[italic]{mathastext}
\usepackage{amsmath}
\usepackage{microtype}
\usepackage[table]{xcolor}
\usepackage{tikz}
\usepackage[most]{tcolorbox}
\usepackage{fontawesome5}
\usepackage{tabularx,booktabs}
\usepackage{enumitem}
\usepackage{fancyhdr}
\usepackage[hidelinks]{hyperref}

% ---------------------------------------------------------------- colours (as in the brochure)
\definecolor{navy}{HTML}{1C3A5E}
\definecolor{pblue}{HTML}{3B78B5}
\definecolor{amber}{HTML}{F2B33D}
\definecolor{deepamber}{HTML}{D9902A}
\definecolor{pale}{HTML}{EEF3F9}
\definecolor{ink}{HTML}{2A2F36}
\definecolor{logonavy}{HTML}{173D63}
\definecolor{logodark}{HTML}{C98A1E}
\definecolor{logoamber}{HTML}{F2B446}
\definecolor{logocream}{HTML}{FEEFD2}
\definecolor{soft}{HTML}{5E6773}
\definecolor{agree}{HTML}{2F7D4F}
\definecolor{differ}{HTML}{B23B3B}
\color{ink}

\setlength{\parindent}{0pt}
\setlength{\parskip}{4pt}
\setlist{nosep,leftmargin=*,itemsep=2pt}
\renewcommand{\arraystretch}{1.15}

\pagestyle{fancy}
\fancyhf{}
\renewcommand{\headrulewidth}{0pt}
\fancyfoot[L]{\footnotesize\color{soft}\textbf{Pollis}\enspace New Methods\enspace·\enspace\nmtheme}
\fancyfoot[R]{\footnotesize\color{soft}\thepage}
\newcommand{\nmtheme}{}
\newcommand{\nmdate}{11 October 2026}

% The Pollis logo, redrawn so it prints sharply at any size
\newcommand{\pollislogo}[1]{%
  \begin{tikzpicture}[x=#1/512,y=-#1/512,baseline=(current bounding box.center)]
    \fill[logonavy] (0,0) rectangle (512,512);
    \fill[logodark] (36,422) -- (104,256) -- (172,422) -- cycle;
    \fill[logoamber] (104,422) -- (203,180) -- (302,422) -- cycle;
    \fill[logocream] (203,422) -- (339,89) -- (476,422) -- cycle;
  \end{tikzpicture}}

% Top of the first page: logo, name; then the kicker, title and a one-line promise
\newcommand{\nmhead}[3]{%
  \renewcommand{\nmtheme}{#1}%
  \thispagestyle{fancy}%
  \begin{minipage}[c]{13mm}\pollislogo{11mm}\end{minipage}%
  \begin{minipage}[c]{\dimexpr\linewidth-13mm\relax}
    {\color{navy}\Large\bfseries Pollis}\enspace{\color{soft}New Methods}\hfill
    {\color{deepamber}\small\bfseries\MakeUppercase{#1}}
  \end{minipage}\par\vspace{7mm}
  {\color{navy}\fontsize{22}{26}\selectfont\bfseries #2}\par\vspace{4pt}
  {\color{soft}\large #3}\par\vspace{5mm}}

\newcommand{\nmsection}[1]{\par\vspace{5pt}{\color{navy}\large\bfseries #1}\par\vspace{1pt}}
\newcommand{\nmsub}[1]{\par\vspace{2pt}{\color{deepamber}\bfseries #1}\par}
\newcommand{\gt}{\,{\color{deepamber}\bfseries ›}\,}
\newcommand{\menupath}[1]{{\bfseries\color{navy}#1}}
\newcommand{\fn}[1]{\texttt{\small #1}}
% Citations: \nmcite{key}{Breiman (2001)} links to the entry \nmref{key}{...} in the References
\newcommand{\nmcite}[2]{\hyperlink{ref:#1}{\color{pblue}#2}}
\newcommand{\nmref}[2]{\raisebox{1.1em}[0pt][0pt]{\hypertarget{ref:#1}{}}#2\par}
\newcommand{\nmdoi}[1]{\href{https://doi.org/#1}{\color{pblue}doi:#1}}
% Where a notebook is in Pollis: its title in Browse Methods, and its file
\newcommand{\browse}[2]{\menupath{Compose \gt New Methods \gt Browse Methods}: \emph{#1} (\fn{#2.ipynb})}

% Key facts at the top: source, missing in Julia, notebook, result
\newcommand{\factsheet}[4]{%
  \begin{tcolorbox}[enhanced, colback=pale, colframe=pale, boxrule=0pt, sharp corners,
    left=7pt, right=7pt, top=5pt, bottom=5pt, fontupper=\small]
  \begin{tabularx}{\linewidth}{@{}>{\leavevmode\bfseries\color{navy}}p{27mm}>{\leavevmode}X@{}}
  Tested against & #1\\
  Missing in Julia & #2\\
  Notebook & #3\\
  Result & #4
  \end{tabularx}
  \end{tcolorbox}}

% Boxes
\newtcolorbox{pbox}[2][]{enhanced, colback=pale, colframe=amber, boxrule=0pt,
  leftrule=3pt, sharp corners, left=7pt, right=7pt, top=6pt, bottom=6pt,
  fontupper=\small\raggedright, before upper={{\color{navy}\normalsize\bfseries #2}\par\smallskip}, #1}
\newtcolorbox{nbox}[2][]{enhanced, colback=navy, colframe=navy, boxrule=0pt,
  sharp corners, left=8pt, right=8pt, top=7pt, bottom=7pt, coltext=white,
  fontupper=\small\raggedright, before upper={{\color{amber}\normalsize\bfseries #2}\par\smallskip}, #1}

% Value tables: published beside Julia
\newenvironment{valuetable}[1]{%
  \par\small\tabularx{\linewidth}{#1}\toprule}{\bottomrule\endtabularx\par}
\newcommand{\ok}{{\color{agree}\faCheck}}
\newcommand{\no}{{\color{differ}\faTimes}}

% End of the document: credit and licence
\newcommand{\nmcredit}{%
  \par\vfill
  \begin{minipage}[b]{0.78\linewidth}\footnotesize\color{soft}
  Methods written and tested with Claude Code (Anthropic) in Pollis, from the published papers and their data.\par
  \textcopyright\ 2026 Antonio Saragga Seabra\enspace·\enspace Licensed under \href{https://creativecommons.org/licenses/by-sa/4.0/}{CC BY-SA 4.0}\enspace·\enspace Draft, \nmdate
  \end{minipage}\hfill
  \begin{minipage}[b]{0.18\linewidth}\raggedleft\pollislogo{9mm}\par\vspace{2pt}{\footnotesize contact@pollis.org}\end{minipage}}


\begin{document}

\nmhead{GLM Inference}
  {Parametric bootstrap and delta-beta influence for binomial GLMs}
  {The pre-Challenger risk analysis of Dalal, Fowlkes and Hoadley (1989), reproduced.}

\factsheet
  {\nmcite{dalal1989}{Dalal, Fowlkes and Hoadley (1989)}, \emph{Journal of the American Statistical Association}: Model (3.2), the bootstrap
   interval at 30\,\textdegree F (p.\,950, Fig.\,5) and the delta beta of each flight (Fig.\,7). The 23 flights of Table 1, from the
   \nmcite{commission1986}{Presidential Commission (1986)} (PollisDatasets \fn{challenger}).}
  {a parametric bootstrap and case-deletion influence statistics for generalized linear models. GLM.jl fits the models only.}
  {\browse{Parametric bootstrap and delta-beta influence for binomial GLMs}{glmextras}}
  {every published value: the fit of Model (3.2), the interval at 30\,\textdegree F, (1.27, 5.95) for ``about (1, 6)'', and
   the delta beta of flight 61-A, 4.22 for ``about four''.}

\nmsection{The problem}
On the morning of 28 January 1986 the temperature at the launch pad was about 31\,\textdegree F, far below any earlier
launch. \nmcite{dalal1989}{Dalal, Fowlkes and Hoadley (1989)} asked what the data available before the launch said about O-ring damage at that
temperature, and how sure one could be. Two tools answer: a \textbf{parametric bootstrap}, for an interval on a prediction
far outside the data, where asymptotic formulas are least trustworthy; and \textbf{delta beta}, to see whether a single
flight drives the fit.

\nmsection{The method}
\nmsub{Parametric bootstrap}
For a binomial GLM with $n$ trials per case, \fn{parametric} takes the fitted probabilities $\hat p$ as true, draws
$y^* \sim \mathrm{Binomial}(n, \hat p)/n$ for every case, refits, and repeats $B$ times \nmcite{efron1979}{Efron (1979)}. The spread of any
refitted quantity, a coefficient or a prediction, estimates its sampling distribution. A sample whose fit fails, as one with
all successes or none can, is skipped and counted.
\nmsub{Delta beta}
Case $i$'s influence on coefficient $j$ is the change when the case is dropped and the model refitted:
$\delta_i = (\hat\beta_j - \hat\beta_{j(-i)})/S$, with $S$ the standard deviation of the $n$ refitted coefficients (p.\,950).
Unlike the one-step approximation of \nmcite{pregibon1981}{Pregibon (1981)}, it refits exactly, which matters when a few cases carry the fit.

\nmsection{What the tests show}
\begin{valuetable}{@{}lXl@{}}
\textbf{Test} & \textbf{What is compared} & \textbf{Agreement}\\\midrule
Model (3.2), p.\,948 & coefficients, standard errors, $G^2$ & every printed digit\\
Fig.\,5, p.\,950 & 90\% interval for damaged O-rings at 30\,\textdegree F & (1.27, 5.95) for ``about (1, 6)''\\
Fig.\,7, p.\,950 & the most influential flight & 61-A, 4.22 s.d.\ for ``about four''\\
Fig.\,7, p.\,950 & the next three & flights 2, 11 and 14, as published\\
\end{valuetable}
The tests are coarse because the paper's are: it prints the interval and the delta beta to no decimals and shows the rest in
figures. They confirm that the methods do what the paper did, not their last digits. The bootstrap is random: its interval
moves in the second decimal with the seed (1,000 samples are drawn).

\nmsection{What the replication adds}
The paper's ``about (1, 6)'' becomes a number: of six field-joint O-rings, between 1.27 and 5.95 expected damaged at
30\,\textdegree F. And the influence of flight 61-A can be asked directly: dropped, its damage at a warm 75\,\textdegree F no
longer softens the fit, the model expects nearly all six O-rings to be damaged, and the interval narrows to the top of the
scale. The notebook's example shows both intervals.

\nmsection{Still open}
The paper also bootstraps a model of primary erosion with temperature and pressure (p.\,953). About a third of its samples
have no maximum likelihood estimate: one coefficient runs off to infinity (separation), and the paper does not say what its
program did with them. Those intervals are not used as tests; the Logistic Regression workflow shows them.

\nmsection{Further considerations}
\begin{itemize}
  \item \textbf{Separation.} The failed samples of the pressure model are a known weakness of maximum likelihood with few
    events. The penalised likelihood of \nmcite{firth1993}{Firth (1993)} would give every sample a finite estimate; it belongs here only once a published
    table tests it.
  \item \textbf{Beyond binomial.} Both functions need only a family's random draw and a refit, so Poisson and gamma GLMs
    follow; each needs its own published test.
  \item \textbf{Upstream.} Both are a few lines over GLM.jl and could be offered to it, or to a GLM diagnostics package,
    with the Challenger values as tests.
\end{itemize}

\nmsection{References}
{\footnotesize
\nmref{dalal1989}{Dalal, S.\,R., Fowlkes, E.\,B. and Hoadley, B. (1989). Risk analysis of the space shuttle: pre-Challenger prediction of failure. \emph{Journal of the American Statistical Association} 84(408), 945--957. \nmdoi{10.1080/01621459.1989.10478858}}
\nmref{efron1979}{Efron, B. (1979). Bootstrap methods: another look at the jackknife. \emph{The Annals of Statistics} 7(1), 1--26.}
\nmref{firth1993}{Firth, D. (1993). Bias reduction of maximum likelihood estimates. \emph{Biometrika} 80(1), 27--38.}
\nmref{pregibon1981}{Pregibon, D. (1981). Logistic regression diagnostics. \emph{The Annals of Statistics} 9(4), 705--724.}
\nmref{commission1986}{Presidential Commission on the Space Shuttle Challenger Accident (1986). \emph{Report to the President}, Vol.\,1, pp.\,129--131: the data.}
}

\nmcredit
\end{document}
